\subsection{Facts about OpenWRT}

\begin{frame}{Facts about OpenWRT}
\begin{center}
\LARGE {\color{fh-ooe-red}OpenWRT}
\href{https://openwrt.org/}{\beamergotobutton{https://openwrt.org/}}
\end{center}

\begin{itemize}
  \item Is an \alert{open source firmware/operating system} for (WiFi) \alert{router}.
  \item Is used on \alert{most consumer routers} as base firmware (often modified by vendor).
  \item Provides a largely \alert{POSIX} \href{https://en.wikipedia.org/wiki/POSIX}{\beamergotobutton{https://en.wikipedia.org/wiki/POSIX}} conform Linux kernel and standard libraries.
  \item Uses the small footprint full POSIX conform \alert{musl libc} \href{https://musl.libc.org/}{\beamergotobutton{https://musl.libc.org/}} instead of the \alert{GNU C library} \href{https://www.gnu.org/software/libc/}{\beamergotobutton{https://www.gnu.org/software/libc/}} as standard C library.
  \begin{description}[]
  \item[$\Rightarrow$] Thus, in principal, all software existing for Linux can be brought to OpenWRT based devices.
  \item[$\Rightarrow$] Limiting factors: Resources of the device
  \begin{itemize}
    \item Memory (RAM)
    \item Memory (FLASH)
    \item System on a Chip (SOC) performance
    \item Small differences in the behaviour of \alert{musl libc} and \alert{GNU libc}
  \end{itemize}
  \end{description}
\end{itemize}
\end{frame}
